
\documentclass{article}
\usepackage{colortbl}
\usepackage{makecell}
\usepackage{multirow}
\usepackage{supertabular}

\begin{document}

\newcounter{utterance}

\centering \large Interaction Transcript for game `hot\_air\_balloon', experiment `air\_balloon\_survival\_en\_negotiation\_easy', episode 5 with Qwen3{-}32B{-}t0.0.
\vspace{24pt}

{ \footnotesize  \setcounter{utterance}{1}
\setlength{\tabcolsep}{0pt}
\begin{supertabular}{c@{$\;$}|p{.15\linewidth}@{}p{.15\linewidth}p{.15\linewidth}p{.15\linewidth}p{.15\linewidth}p{.15\linewidth}}
    \# & $\;$A & \multicolumn{4}{c}{Game Master} & $\;\:$B\\
    \hline

    \theutterance \stepcounter{utterance}  
    & & \multicolumn{4}{p{0.6\linewidth}}{
        \cellcolor[rgb]{0.9,0.9,0.9}{
            \makecell[{{p{\linewidth}}}]{
                \texttt{\tiny{[P1$\langle$GM]}}
                \texttt{You are participating in a collaborative negotiation game.} \\
\\ 
\texttt{Together with another participant, you must agree on a single set of items that will be kept. Each of you has your own view of how much each item matters to you (importance). You do not know how the other participant values the items. Additionally, you are given the effort each item demands.} \\
\texttt{You may only agree on a set if the total effort of the selected items does not exceed a shared limit:} \\
\\ 
\texttt{LIMIT = 3170} \\
\\ 
\texttt{Here are the individual item effort values:} \\
\\ 
\texttt{Item effort = \{"C08": 424, "C55": 118, "A43": 846, "A67": 907, "B42": 809, "A62": 41, "B19": 193, "C43": 246, "C73": 805, "C10": 601, "B36": 432, "A18": 166, "C61": 119, "C67": 462, "C42": 172\}} \\
\\ 
\texttt{Here is your personal view on the importance of each item:} \\
\\ 
\texttt{Item importance values = \{"C08": 138, "C55": 583, "A43": 868, "A67": 822, "B42": 783, "A62": 65, "B19": 262, "C43": 121, "C73": 508, "C10": 780, "B36": 461, "A18": 484, "C61": 668, "C67": 389, "C42": 808\}} \\
\\ 
\texttt{Goal:} \\
\\ 
\texttt{Your goal is to negotiate a shared set of items that benefits you as much as possible (i.e., maximizes total importance to YOU), while staying within the LIMIT. You are not required to make a PROPOSAL in every message {-} you can simply negotiate as well. All tactics are allowed!} \\
\\ 
\texttt{Interaction Protocol:} \\
\\ 
\texttt{You may only use the following structured formats in a message:} \\
\\ 
\texttt{PROPOSAL: \{'A', 'B', 'C', …\}} \\
\texttt{Propose keeping exactly those items.} \\
\texttt{REFUSE: \{'A', 'B', 'C', …\}} \\
\texttt{Explicitly reject opponent's proposal.} \\
\texttt{ARGUMENT: \{'...'\}} \\
\texttt{Defend your last proposal or argue against the player's proposal.} \\
\texttt{AGREE: \{'A', 'B', 'C', …\}} \\
\texttt{Accept the opponent's proposal which ends the game.} \\
\texttt{STRATEGIC REASONING: \{'...'\}} \\
\texttt{	Describe your strategic reasoning or anticipation explaining your choice of action. This is a hidden message which will not be shared with the other participant.} \\
\\ 
\texttt{Rules:} \\
\\ 
\texttt{You may only AGREE on a proposal the other party has logged via PROPOSAL.} \\
\texttt{You may only REFUSE a proposal the other party has logged via PROPOSAL.} \\
\texttt{Total effort of any PROPOSAL or AGREE set must be ≤ LIMIT.} \\
\texttt{Do NOT reveal your hidden importance scores.} \\
\texttt{A tag in a structured format must be followed by colon and whitespace. The argument must be a python set containing 0 or more strings.} \\
\texttt{So, it must be of the form TAG: \{...\}} \\
\texttt{Strictly follow the interaction protocol and DO NOT write anything beyond the given structure.} \\
\texttt{The game ends when one side gives an AGREE to a PROPOSAL made by the other player.} \\
\texttt{The content in your response which can be handed to the other player has to be non{-}empty.} \\
\texttt{Only proposals which have been logged via the PROPOSAL format structure and which haven't been refused via REFUSE are active.} \\
\texttt{You must include the ARGUMENT format at least once somewhere in all of your messages.} \\
\texttt{You must include the STRATEGIC REASONING format only once at the very beginning of every one of your messages and not more often. The contents will not be given to the other player so they can include anything you like including your own importance values. Here you should reason step by step to come up with you next move.} \\
            }
        }
    }
    & \\ \\

    \theutterance \stepcounter{utterance}  
    & \multicolumn{4}{p{0.6\linewidth}}{
        \cellcolor[rgb]{0.9,0.9,0.9}{
            \makecell[{{p{\linewidth}}}]{
                \texttt{\tiny{[P1$\rangle$GM]}}
                \texttt{<think>} \\
\texttt{Okay, let's see. I need to figure out the best set of items to propose that maximizes my total importance without exceeding the effort limit of 3170. First, I should calculate the total effort of all items to see if they can all fit. Let me add them up.} \\
\\ 
\texttt{Looking at the item efforts: C08 is 424, C55 is 118, A43 is 846, A67 is 907, B42 is 809, A62 is 41, B19 is 193, C43 is 246, C73 is 805, C10 is 601, B36 is 432, A18 is 166, C61 is 119, C67 is 462, C42 is 172. Let me sum these up step by step.} \\
\\ 
\texttt{Starting with 0. Adding 424 (C08) gives 424. Then +118 (C55) is 542. +846 (A43) is 1388. +907 (A67) is 2295. +809 (B42) is 3104. +41 (A62) is 3145. +193 (B19) is 3338. Oh, that's already over 3170. So including all items is impossible. So I need to exclude some items.} \\
\\ 
\texttt{The problem is to select a subset with total effort ≤3170 and maximize my total importance. Since I don't know the other person's preferences, I need to prioritize items that have high importance for me but maybe lower effort, or find a balance.} \\
\\ 
\texttt{Let me list the items with their importance and effort:} \\
\\ 
\texttt{{-} A43: 868 importance, 846 effort} \\
\texttt{{-} A67: 822, 907} \\
\texttt{{-} B42: 783, 809} \\
\texttt{{-} C10: 780, 601} \\
\texttt{{-} A18: 484, 166} \\
\texttt{{-} B36} \\
            }
        }
    }
    & & \\ \\

    \theutterance \stepcounter{utterance}  
    & & \multicolumn{4}{p{0.6\linewidth}}{
        \cellcolor[rgb]{0.9,0.9,0.9}{
            \makecell[{{p{\linewidth}}}]{
                \texttt{\tiny{[P1$\langle$GM]}}
                \texttt{Your response did not start with the proper strategic reasoning tag at the very beginning of your response. The very first structured format must be of the form STRATEGIC REASONING: \{...\}. Try again.} \\
            }
        }
    }
    & \\ \\

    \theutterance \stepcounter{utterance}  
    & \multicolumn{4}{p{0.6\linewidth}}{
        \cellcolor[rgb]{0.9,0.9,0.9}{
            \makecell[{{p{\linewidth}}}]{
                \texttt{\tiny{[P1$\rangle$GM]}}
                \texttt{<think>} \\
\texttt{Okay, let me start over. The user wants me to make sure the first part is the STRATEGIC REASONING. So first, I need to outline my plan.} \\
\\ 
\texttt{I need to maximize my total importance while staying under the effort limit. Let me list all items with their importance and effort. Then, maybe sort them by importance per unit effort to prioritize high{-}value items. But since the other person's preferences are unknown, I should focus on my own high{-}importance items that don't add too much effort.} \\
\\ 
\texttt{Looking at my importance values: A43 (868), A67 (822), B42 (783), C10 (780), A18 (484), B36 (461), C61 (668), C55 (583), C08 (138), etc. The top ones are A43, A67, B42, C10, C61, A18, B36, C55.} \\
\\ 
\texttt{Now, check their efforts. A43 is 846, A67 is 907, B42 is 809, C10 is 601, C61 is 119, A18 is 166, B36 is 432, C55 is 118. } \\
\\ 
\texttt{If I take A43 (846), A67 (907), B42 (809), that's 846+907+809 = 2562. Then adding C10 (601) would make 3163, which is under 3170. That's total effort 3163. Let's check the total importance: 868 + 822 + 783 + 780 = 3253. Then, maybe add C61 (119 effort, 668 importance) but 3163 +119=3282 which is over. So can't add that. Alternatively, replace some items. Maybe remove A67 (907) and add C61 (119) and others. Let's see: 84} \\
            }
        }
    }
    & & \\ \\

    \theutterance \stepcounter{utterance}  
    & & \multicolumn{4}{p{0.6\linewidth}}{
        \cellcolor[rgb]{0.9,0.9,0.9}{
            \makecell[{{p{\linewidth}}}]{
                \texttt{\tiny{[P1$\langle$GM]}}
                \texttt{Your response did not start with the proper strategic reasoning tag at the very beginning of your response. The very first structured format must be of the form STRATEGIC REASONING: \{...\}. Try again.} \\
            }
        }
    }
    & \\ \\

    \theutterance \stepcounter{utterance}  
    & \multicolumn{4}{p{0.6\linewidth}}{
        \cellcolor[rgb]{0.9,0.9,0.9}{
            \makecell[{{p{\linewidth}}}]{
                \texttt{\tiny{[P1$\rangle$GM]}}
                \texttt{<think>} \\
\texttt{Okay, I need to make sure I start with the STRATEGIC REASONING tag. Let me outline my plan step by step.} \\
\\ 
\texttt{First, I'll calculate the total effort of all items to see if they can all fit. Adding them up: C08 (424) + C55 (118) + A43 (846) + A67 (907) + B42 (809) + A62 (41) + B19 (193) + C43 (246) + C73 (805) + C10 (601) + B36 (432) + A18 (166) + C61 (119) + C67 (462) + C42 (172). Let me add these step by step. 424 + 118 = 542. +846 = 1388. +907 = 2295. +809 = 3104. +41 = 3145. +193 = 3338. That's already over 3170. So I need to exclude some items.} \\
\\ 
\texttt{Next, I'll prioritize items with high importance for me. My top importance items are A43 (868), A67 (822), B42 (783), C10 (780), C61 (668), A18 (484), B36 (461), C55 (583), C08 (138), etc. Now, check their efforts. A43 is 846, A67 is 907, B42 is 809, C10 is 601, C61 is 119, A18 is 166, B36 is 432, C55 is 118.} \\
\\ 
\texttt{If I take A43, A67, B42, and C10, their total effort is 846 + 907 + 809 + 601 = 3163,} \\
            }
        }
    }
    & & \\ \\

    \theutterance \stepcounter{utterance}  
    & & & \multicolumn{2}{p{0.3\linewidth}}{
        \cellcolor[rgb]{0.9,0.9,0.9}{
            \makecell[{{p{\linewidth}}}]{
                \texttt{\tiny{[GM$|$GM]}}
                \texttt{abort} \\
            }
        }
    }
    & & \\ \\

    \theutterance \stepcounter{utterance}  
    & & & \multicolumn{2}{p{0.3\linewidth}}{
        \cellcolor[rgb]{0.9,0.9,0.9}{
            \makecell[{{p{\linewidth}}}]{
                \texttt{\tiny{[GM$|$GM]}}
                \texttt{end game} \\
            }
        }
    }
    & & \\ \\

\end{supertabular}
}

\end{document}
